\lysh{14562}{4}

\begin{alist}
\item 
\begin{rlist}
\item Η συνάρτηση ορίζεται για τις τιμές του $x \in \mathbb{R}$, για τις οποίες ισχύει $x^2 - 3x + 2 \ne 0 \quad (1)$.
Το τριώνυμο $x^2 - 3x + 2$ έχει $\alpha = 1, \beta = -3, \gamma = 2$ και $\Delta = \beta^2 - 4\alpha\gamma = (-3)^2 - 4 \cdot 1 \cdot 2 = 1 > 0$.
Οι ρίζες του είναι $x_1 = \dfrac{-\beta - \sqrt{\Delta}}{2\alpha} = 1$ και $x_2 = \dfrac{-\beta + \sqrt{\Delta}}{2\alpha} = 2$.
Άρα η $(1)$ ισχύει για $x \ne 1$ και $x \ne 2$. Συνεπώς το πεδίο ορισμού της συνάρτησης είναι:
$A = \mathbb{R} - \{1, 2\}$.

\item Είναι $f(x) = \dfrac{x^2 - x}{x^2 - 3x + 2} = \dfrac{x(x - 1)}{(x - 1)(x - 2)} = \dfrac{x}{x - 2}$, $x \in \mathbb{R} - \{1, 2\}$.
\end{rlist}

\item Οι τετμημένες των κοινών σημείων, αν αυτά υπάρχουν, θα είναι οι λύσεις της εξίσωσης:
$|f(x)| = 1$, δηλαδή $\left|\dfrac{x}{x - 2}\right| = 1$, οπότε $\dfrac{x}{x - 2} = 1$ ή $\dfrac{x}{x - 2} = -1$, δηλαδή
$x = x - 2$ ή $x = -(x - 2)$ και τελικά
$0x = -2$, που είναι αδύνατη ή $x = 1$ που δεν είναι δεκτή λύση (αφού $1 \notin A$).
Συνεπώς οι ευθείες $y = 1$ και $y = -1$ δεν έχουν κοινά σημεία με τη γραφική παράσταση της $f(x)$.
\end{alist}

